\chapter{Adic coefficients, Weil structures, and Frobenius}
\label{ch:adic-weil}

Fix a prime \(\ell\) invertible on every scheme under consideration and put
\[
 \Lambda_n=\mathbf Z/\ell^n\mathbf Z,\qquad
 \Lambda_\bullet=(\Lambda_n)_{n\geq1}.
\]
An arbitrary inverse system of constructible sheaves is not an
\(\ell\)-adic sheaf: inverse limit is not exact, and pro-zero errors must be
removed.  The normalized Artin--Rees formalism does exactly this.

\section{Normalized adic systems}

\begin{definition}[Strict adic system]
\label{def:strict-adic-system}
A strict constructible \(\mathbf Z_\ell\)-system is a projective system
\(\mathcal F=(\mathcal F_n)\) of constructible \(\Lambda_n\)-sheaves such that
\[
 \Lambda_n\otimes_{\Lambda_{n+1}}\mathcal F_{n+1}
 \xrightarrow{\sim}\mathcal F_n.
\]
For complexes the normalization condition is the derived isomorphism
\[
 \Lambda_n\otimes^L_{\Lambda_{n+1}}K_{n+1}
 \xrightarrow{\sim}K_n.
\]
\end{definition}

A projective system \(\mathcal M\) is \emph{AR-null} if, locally on a finite
constructible stratification, for every \(n\) some transition
\(\mathcal M_{n+r}\to\mathcal M_n\) is zero.  A complex is AR-null when all its
cohomology systems are AR-null.  It is AR-adic when, modulo an AR-null
complex, it has a bounded constructible derived-normalized representative.

\begin{definition}[The normalized adic category]
\label{def:normalized-adic}
Let \(D^b_{\mathrm{AR}}(X,\Lambda_\bullet)\) be the category of bounded
AR-adic complexes and \(D^b_{\mathrm{null}}\) its thick subcategory of
AR-null complexes.  Define
\[
 D^b_c(X,\mathbf Z_\ell)
 :=
 D^b_{\mathrm{AR}}(X,\Lambda_\bullet)/
 D^b_{\mathrm{null}}(X,\Lambda_\bullet).
\]
\end{definition}

\begin{proposition}[Normalization and reduction]
Every object has a normalized representative, unique up to an AR-null cone.
Reduction defines conservative functors
\[
 (-)_n:D^b_c(X,\mathbf Z_\ell)\longrightarrow D^b_c(X,\Lambda_n).
\]
Thus a map between normalized objects is an isomorphism if and only if all its
finite-level reductions are isomorphisms.
\end{proposition}

\begin{proof}[Proof capsule]
Replace an AR-adic representative by its derived \(\ell\)-adic completion
formed from compatible finite-level truncations.  Artin--Rees implies that
the comparison cone is AR-null, and the completed system satisfies derived
normalization.  If every reduction of a normalized cone vanishes, its
cohomology pro-systems are AR-null; hence the cone is zero in the quotient.
See \cite{KiehlWeissauer}.
\end{proof}

\begin{definition}[\(\mathbf Q_\ell\)-complexes]
\label{def:adic-derived}
Set
\[
 D^b_c(X,\mathbf Q_\ell)
 :=D^b_c(X,\mathbf Z_\ell)\otimes_{\mathbf Z_\ell}\mathbf Q_\ell,
\]
meaning that lattices are retained as objects and Hom groups are tensored
with \(\mathbf Q_\ell\).  Two commensurable stable lattices therefore define
the same rational object.
\end{definition}

\section{The adic six operations}

\begin{theorem}[Adic six operations]
\label{thm:adic-six-ops}
For a separated finite-type morphism \(f:X\to Y\), the finite-level operations
induce
\[
 Rf_*,\quad Rf_!,\quad f^*,\quad Rf^!,\quad
 -\otimes^L-,\quad R\mathcal Hom(-,-)
\]
on the constructible \(\mathbf Z_\ell\)- and \(\mathbf Q_\ell\)-categories,
with the appropriate finite-Tor-amplitude hypotheses for internal Hom,
\(Rf^!\), and duality.  They satisfy adjunction, composition, localization,
base change, projection formula, and Kunneth.  Reduction modulo \(\ell^n\)
commutes canonically with each operation.
\end{theorem}

\begin{proof}[Proof capsule]
Apply each operation to normalized representatives level by level and
normalize.  For \(\star\in\{*,!\}\), finite-level base change in the
coefficient ring and projection formula give
\[
 \Lambda_n\otimes^L_{\Lambda_{n+1}}Rf_\star K_{n+1}
 \simeq Rf_\star
  (\Lambda_n\otimes^L_{\Lambda_{n+1}}K_{n+1})
 \simeq Rf_\star K_n.
\]
Uniform cohomological-dimension bounds keep totalizations bounded.
Artin--Rees shows that the operations preserve AR-null cones.  Right adjoints
descend by adjunction.  Every compatibility is an isomorphism after all
finite reductions, and therefore adicly by conservativity.  Finally invert
\(\ell\).
\end{proof}

\begin{theorem}[Finiteness and perfectness]
\label{thm:adic-perfectness}
Let \(X\) be separated of finite type over a finite field and
\(K\in D^b_c(X,\mathbf Q_\ell)\).  Then
\[
 R\Gamma_c(X_{\bar k},K)
\]
is perfect over \(\mathbf Q_\ell\); in particular its cohomology is
finite-dimensional and bounded.  For a constructible
\(\mathbf Z_\ell\)-lattice of finite Tor amplitude, compactly supported
cohomology is a perfect \(\mathbf Z_\ell\)-complex.
\end{theorem}

\begin{proof}[Proof capsule]
Finite-coefficient finiteness gives finite \(\Lambda_n\)-cohomology in a
range independent of \(n\).  Normalization and the derived
Mittag--Leffler triangle identify the derived inverse limit with adic
cohomology.  Topological Nakayama gives finite generation, while devissage on
a finite stratification and the uniform dimension bound give finite Tor
amplitude.  Over the discrete valuation ring \(\mathbf Z_\ell\), such a
bounded complex is perfect; invert \(\ell\).
\end{proof}

\section{The Frobenius dictionary}

Let \(k_0=\mathbf F_q\), \(k=\overline{\mathbf F}_q\), and
\(X=X_0\times_{k_0}k\).  The following objects are distinct.
\begin{enumerate}
\item The arithmetic Galois element
\(\sigma\in\operatorname{Gal}(k/k_0)\), \(a\mapsto a^q\).
\item Its inverse \(\sigma^{-1}\), called the geometric Galois Frobenius.
\item The \(q\)-power Frobenius morphism
\(\operatorname{Fr}_{X_0}:X_0\to X_0\), whose pullback on functions is
\(f\mapsto f^q\).
\item Its base change
\[
 F_X:X\longrightarrow X.
\]
On geometric coordinates \(F_X\) sends \(x\) to \(x^q\); it is a finite
universal homeomorphism and \(dF_X=0\).
\item For a closed point \(x\in|X_0|\), local geometric Frobenius \(F_x\) in
the decomposition group.
\end{enumerate}

The scheme map in item 4 acts covariantly on points but cohomology is
contravariant.  The standard descent comparison identifies the operator
induced by \(F_X^*\), together with the Weil descent isomorphism, with the
action of the geometric Galois element \(\sigma^{-1}\).  Thus the same symbol
\(F\) may consistently denote the \(q\)-power scheme correspondence and the
resulting \emph{geometric} Frobenius operator on stalks and cohomology.

\begin{convention}[title={Geometric Frobenius}]
\label{conv:geometric-frobenius}
Every unqualified \(F\) acting on a stalk or cohomology group is geometric
Frobenius.  If \(y\in X_0(\mathbf F_{q^m})\) has underlying closed point
\(x\) of degree \(d\mid m\), then the global \(F^m\) on the fixed-point stalk
is
\[
 F_y^m=F_x^{\,m/d}.
\]
\end{convention}

\section{Tate twists}

\begin{definition}[Tate objects]
\label{def:tate-twist}
Define
\[
 \mathbf Z_\ell(1)=\varprojlim_n\mu_{\ell^n},\qquad
 \mathbf Q_\ell(1)=\mathbf Z_\ell(1)\otimes\mathbf Q_\ell,
\]
and define \(K(r)\) by tensor powers for \(r\geq0\) and dual tensor powers for
\(r<0\).
\end{definition}

\begin{proposition}[Frobenius on twists]
\label{prop:frob-tate}
Geometric Frobenius acts on \(\mathbf Q_\ell(1)\) by \(q^{-1}\), hence on
\(\mathbf Q_\ell(r)\) by \(q^{-r}\).  Twisting an eigenvalue \(\alpha\) by
\((r)\) replaces it by \(q^{-r}\alpha\).
\end{proposition}

\begin{proof}
Arithmetic Galois Frobenius acts on roots of unity by
\(\zeta\mapsto\zeta^q\); geometric Galois Frobenius acts by the inverse of
\(q\) modulo every \(\ell^n\).  Pass to the inverse limit, then use tensor
powers and duals.
\end{proof}

\section{Weil complexes}

\begin{definition}[Weil complex]
\label{def:weil-complex}
A Weil complex on \(X_0\) is a pair \((K,\phi)\), with
\(K\in D^b_c(X,\mathbf Q_\ell)\) and an isomorphism
\[
 \phi:F_X^*K\xrightarrow{\sim}K.
\]
Its \(m\)-th iterate is the ordered composite
\(\phi_m:(F_X^m)^*K\to K\).
\end{definition}

For an object descended from \(X_0\), the canonical descent structure provides
\(\phi\).  At a fixed geometric point, \(\phi_m\) is the geometric Frobenius
operator used in the trace formula.

\begin{proposition}[Weil structures and six operations]
\label{prop:weil-six-ops}
The six operations carry canonical Weil structures.  For example,
\[
 F_Y^*Rf_!K\simeq Rf_!F_X^*K
 \xrightarrow{Rf_!(\phi)}Rf_!K
\]
is the Weil structure on \(Rf_!K\).  The analogous structures on pullback,
exceptional pullback, tensor, internal Hom, and duality respect all adjunction
and base-change morphisms.
\end{proposition}

\begin{proof}
The square formed by \(f,F_X,F_Y\) commutes, but it need not be Cartesian.
Let \(X\to X\times_{Y,F_Y}Y\) be the relative Frobenius.  It is a
universal homeomorphism.  Apply base change to the genuinely Cartesian
square, then use topological invariance of the etale topos along relative
Frobenius; this gives the displayed exchange isomorphism.  Compose with
\(Rf_!(\phi)\).  Iteration and coherence hold at every finite level and
therefore in the adic quotient.
\end{proof}

\section{Formalization checklist}

\begin{enumerate}
\item Implement coefficient towers, derived normalization, AR-null complexes,
and the Verdier quotient; never identify inverse limit with an exact functor.
\item Prove conservativity of all finite reductions before transporting
coherence to the adic category.
\item Keep integral lattices separate from rational objects and state exactly
where finite Tor amplitude is required.
\item Give distinct definitions for arithmetic and geometric Galois
Frobenius, the \(q\)-power scheme morphism, its induced cohomological operator,
and local closed-point Frobenius.
\item Export the comparison between \(F_X^*\) and geometric Galois action and
the conversion \(F_y^m=F_x^{m/\deg x}\).
\item Test
\[
 H^0(\operatorname{Spec}k,\mathbf Q_\ell),\quad
 \mathbf Q_\ell(1),\quad
 H^2(\mathbf P^1_k,\mathbf Q_\ell)=\mathbf Q_\ell(-1);
\]
the eigenvalues must be \(1,q^{-1},q\).
\end{enumerate}


